% !TeX program = xelatex
% UTF-8; standalone XeLaTeX/XePersian document.
% Compile from the repository root or this file's directory.
% Define \FeverInNotebook before \input to include in a XePersian document.
% Source: Google Drive Notebook/fever, handwritten page order.
% pic01.jpg = 20261012_174550722.jpg; 1XCoCaEWmEmO1Gq__8Pnsf2zK3XthrOJB
% pic02.jpg = 20261012_174556513.jpg; 1UzscDPZZSN6cChhGcOE2y2Rg5DRzcUXB
% pic03.jpg = 20261012_174600923.jpg; 1ELBDxQROv4-kvzlFbJp0tOjv7HJuqGeh
% pic04.jpg = 20261012_174608710.jpg; 1ybjwTOde2nE1gdoMG2IOKOD2h4lEYCVq
% pic05.jpg = 20261012_174613844.jpg; 1EHKH9SN1OLsDKdGAqhg7jhyqXxbntP--
% pic06.jpg = 20261012_174618272.jpg; 13JV9PCG45kioppZ15rEx8vTrPKDsg7wN
\ifdefined\FeverInNotebook
  \let\FeverFinish\relax
\else
  \def\FeverFinish{\end{document}}
  \documentclass[12pt,a4paper]{article}
  \usepackage[margin=22mm]{geometry}
  \usepackage{amsmath,amssymb}
  \usepackage{graphicx}
  \usepackage{xepersian}
  \IfFontExistsTF{Amiri}{\settextfont{Amiri}}{\settextfont{DejaVu Sans}}
  \setlatintextfont{DejaVu Serif}
  \setlength{\parindent}{0pt}
  \setlength{\parskip}{5pt}
  \setlength{\emergencystretch}{3em}
  \begin{document}
\fi
\providecommand{\FeverUnclear}[1]{\textbf{[خوانش نامطمئن: #1]}}
\providecommand{\FeverMath}[1]{\lr{\(#1\)}}
\providecommand{\FeverPage}[2]{\clearpage\section*{تصویر #1: \lr{\texttt{pic#2.jpg}}}}
\providecommand{\FeverHeader}[1]{\textbf{سربرگ و حاشیه:} تب ـ ۲۱۹؛
  \lr{Subject:} (خالی)؛ \lr{Date: / /}؛ شمارهٔ برگه: #1؛ نشان پایین برگه: \lr{YASHA}.\par}

\begin{center}
{\Large\bfseries تب ـ \lr{Fever}}\\
\lr{Nelson 22}\\
۲۱۹
\end{center}
\textbf{یادداشت رونویسی:} متن شش تصویر، شامل سربرگ‌ها و نوشته‌های حاشیه، به ترتیب برگه‌ها آمده است. رابطهٔ فلش‌ها در قالب جمله یا فهرست حفظ شده است. خوانش‌های نامطمئن مشخص شده‌اند و گزاره‌های پزشکی منبع اصلاح نشده‌اند. در علامت‌های مقایسهٔ مبهم، ترتیب افقی نمادها حفظ شده است. تصاویر کامل منبع در پیوست همین سند قرار دارند.

\FeverPage{۱}{01}
\textbf{سربرگ و حاشیه:} تب ـ ۲۱۹؛ \lr{Subject: Nelson 22}؛ \lr{Date: / /}؛ شمارهٔ برگه: ۱؛ \lr{YASHA}.

\subsection*{تعریف و دمای طبیعی}
\begin{itemize}
\item معیار تب: \FeverMath{38\leq} رکتال.
\item تب \FeverMath{40<} درجه، هایپرپایرکسی نام دارد.
\item به طور طبیعی دمای بدن بین \lr{36.6--37.9} متغیر است.
\item کمترین ابتدای صبح، بیشترین آخر شب.
\end{itemize}
\textbf{یادداشت نمادها:} در سطر اول عدد ۳۸ سمت چپ علامت دوخطی و در سطر دوم عدد ۴۰ سمت چپ علامت «<» دیده می‌شود.
\FeverUnclear{رقم اعشار حد پایین دمای طبیعی شبیه ۶ است؛ ۲ نیز خوانش محتمل است.}

\subsection*{مکانیسم‌های ایجاد تب}
\begin{itemize}
\item پایروژن‌ها.
\item نقص در از دست دادن حرارت.
\item بیشتر شدن یا بالاتر رفتن حرارت از از دست دادن آن‌ها.
\end{itemize}
\FeverUnclear{سطر سوم به صورت «بیشتر شدن یا بالاتر رفتن حرارت از از دست دادن آن‌ها» خوانده می‌شود.}

\subsection*{پایروژن‌ها}
\textbf{اندوژن:}
\begin{latin}
IL-1, IL-6, TNF-\(\alpha\), IFN-\(\beta\), IFN-\(\gamma\)\\
PG-E\(_2\)
\end{latin}
\textbf{اصلاح دست‌نویس:} در دو اختصار آخر، \lr{IFN} بالای اختصار خط‌خورده نوشته شده است؛ حرف یونانی آخر \FeverUnclear{شبیه \FeverMath{\gamma} است}.

ایجاد تب از مسیر تولید پایروژن‌ها: عفونت، داروها، بدخیمی، پروسهٔ التهابی.

موادی که پایروژن‌ها را تحریک می‌کنند:
کمپلکس \lr{Ag-Ab} در حضور اجزأ کمپلمان، محصولات لنفوسیت‌ها، اسیدهای صفراوی، متابولیت‌های استروئیدهای آندروژنیک.
\FeverUnclear{«اسیدهای صفراوی» و «متابولیت‌های استروئیدهای آندروژنیک» خوانش محتمل دو عبارت انتهای این فهرست هستند.}

\textbf{اگزوژن:}
\begin{itemize}
\item داروها، میکروب‌ها.
\item اندوتوکسین به طور مستقیم ترمورگولیشن را تغییر می‌دهد.
\end{itemize}

\FeverPage{۲}{02}
\FeverHeader{۲}
\subsection*{داروهایی که شایع تب می‌دهند}
ونکومایسین، آمفوتریسین \lr{B}، آلوپورینول.

\subsection*{شرایطی که تولید حرارت از دفع آن بیشتر می‌شود}
\begin{itemize}
\item مسمومیت با سالیسیلات.
\item هایپرترمی بدخیم.
\end{itemize}

\subsection*{شرایطی که در آن \lr{defective heat loss} رخ می‌دهد}
\begin{itemize}
\item دیسپلازی اکتودرمال.
\item \lr{Severe heat exposure}.
\end{itemize}
\FeverUnclear{واژهٔ انگلیسی نخست در سطر دوم شبیه «Sever» نوشته شده است.}

\subsection*{\lr{Fever Etiology}}
\begin{latin}
\begin{itemize}
\item Infectious
\item Inflammatory
\item Neoplastic
\item Miscellaneous
\end{itemize}
\end{latin}

در کودک \lr{Neurologically intact}، تب بالاتر از ۴۲ درجه نمی‌رود.
\begin{itemize}
\item مگر: تیروتوکسیکوز، بدخیمی.
\end{itemize}

تب در عفونت‌های ویروسی معمولاً به آهستگی ظرف چند روز کاهش می‌یابد.

تب در عفونت‌های باکتریال با \lr{ABT} به سرعت رفع می‌شود.

\FeverPage{۳}{03}
\FeverHeader{۳}
\subsection*{الگوهای تب}
\begin{itemize}
\item تب \lr{Intermittent}: ریتم سیرکادین \lr{Exaggerated} است.
  \FeverUnclear{املای انگلیسی منبع شبیه «Exagerated» است.}
\item تب \lr{Septic} یا \lr{hectic}: \lr{Extremely wide fluctuation}.
\item تب \lr{Sustained}: تب \lr{persistent} که در روز کمتر از ۰٫۵ درجه تغییر می‌کند.
\item تب \lr{Remittent}: روزانه بیش از ۰٫۵ درجه تغییر می‌کند.
\item تب \lr{Relapsing}: اپیزودهای تب با فواصل بدون علامت.
\item تب \lr{Tertian}: روز ۱ و ۳؛ مالاریا (ویواکس).
\item تب \lr{Quartan}: روز ۱ و ۴؛ مالاریا (مالاریه).
\end{itemize}

بیماری با تب \lr{relapsing} باید از عود بیماری افتراق داده شود.

\subsection*{تب \lr{Biphasic}}
یک بیماری که دو دورهٔ تب جداگانه دارد (تب \lr{Camelback}).
\begin{itemize}
\item پولیومیلیت.
\item لپتوسپیروز، تب دنگی، تب زرد.
\end{itemize}

\subsection*{تب \lr{periodic}}
سندرم‌هایی که پریودیسیتی رگولار دارند:
\begin{itemize}
\item نوتروپنی \lr{cyclic}، آفتوز استوماتایتیس، فارنژیت، آدنایتیس.
\item \lr{FMF} (تب مدیترانه‌ای خانوادگی).
\item تب هایبرنین (سندرم دوره‌ای مرتبط با رسپتور \lr{TNF}).
\item سندرم هایپر \lr{IgD}.
\end{itemize}
\FeverUnclear{املای واژهٔ پس از «تب» در مورد TNF به صورت «هایبرنین» خوانده می‌شود.}

\subsection*{تب \lr{Double quotidian}}
روزی ۲ بار اوج می‌گیرد؛ همراه با \lr{JIA}.

\FeverPage{۴}{04}
\FeverHeader{۴}
\subsection*{شدت و زمان تب}
طی یک اپیزود، یک \lr{Spike} تب با بیماری عفونی همراه نیست.
\begin{itemize}
\item بیشتر ترنسفیوژن خون یا دارو.
\end{itemize}
\FeverUnclear{در سطر اول، «اپیزود» بالای سطر اضافه شده و واژهٔ زیر آن خط خورده است؛ جمله با حفظ این اصلاح خوانده شده است.}

تب \FeverMath{41<} درجه معمولاً علل غیرعفونی دارد.

\subsection*{علل تب با نماد \FeverMath{41<}}
\begin{itemize}
\item تب سنترال (آسیب هیپوتالاموس، آسیب نخاع).
\item هایپرترمی بدخیم.
\item \lr{NMS}.
\item تب دارویی.
\item \lr{Heat stroke}.
\end{itemize}
\textbf{یادداشت نماد:} عدد ۴۱ و «<» عین ترتیب افقی منبع آمده‌اند.

\subsection*{تظاهرات بالینی}
\begin{itemize}
\item کودک تب‌دار \lr{pale} است و بی‌اشتها.
\item هر درجه تب، \lr{HR} را ۱۰ تا بالا می‌برد؛ در کودک بالای ۲ ماه.
\item \lr{Relative Tachy}: تاکی‌کاردی با درجهٔ تب متناسب نیست؛ علت معمولاً غیرعفونی یا زمانی که توکسین عامل است.
\item \lr{Relative Brady}: بروسلوز، تیفوئید فیور، لپتوسپیروز، دارو.
\item نقص هدایتی قلب: تب روماتیسمی، لایم، میوکاردیت، اندوکاردیت.
\end{itemize}
\FeverUnclear{در عبارت علت تاکی‌کاردی، واژهٔ پیش از «عامل است» شبیه «توکسین» است.}

\FeverPage{۵}{05}
\FeverHeader{۵}
\subsection*{\lr{Evaluation of Acute Fever}}
\begin{latin}
\begin{itemize}
\item H/x: Onset
\item Ph/E: focus on localizing symptoms
\item Lab: CBC, B/C, CRP, procalcitonin, ESR
\item CSF: cell count, culture, glucose, Gram stain, NAAT (for HSV), protein
\item Stool: Calprotectin, Culture, NAAT for enteric pathogens
\end{itemize}
\end{latin}

عفونت محتمل‌ترین علت تب حاد است.

کودک با فارنژیت: \lr{PCR} استرپ گروه \lr{A} و کشت.

تب به خودی خود به سایر ارگان‌ها آسیب نمی‌زند مگر در شرایط خاص:
استاتوس اپیلپتیکوس و \lr{heat stroke}.
\FeverUnclear{عبارت «به خودی خود» و فعل «نمی‌زند» در این سطر با اتصال حروف خوانده شده‌اند.}

درمان علامتی تب در یک بیماری خودمحدودشونده صرفاً با هدف پایین آوردن دمای بدن در کودک سالم ضروری نیست.

تب به دلیل افزایش مصرف \FeverMath{O_2} و همچنین افزایش کارکرد قلب می‌تواند منجر به دکامپنسیشن نارسایی قلبی شود:
\begin{itemize}
\item افراد با بیماری زمینه‌ای قلب.
\item افراد با آنمی مزمن (\lr{SSD}).
\end{itemize}

\FeverPage{۶}{06}
\FeverHeader{۶}
\subsection*{تب، تشنج و درمان}
کودک از ۶ ماهگی تا ۵ سالگی در خطر افزایش یا تشدید \lr{simple febrile seizure}.
\FeverUnclear{واژهٔ متصل به «یا» شبیه «تشدید» است؛ عبارت آغاز سطر «در خطر افزایش» خوانده می‌شود.}

کودک با ایدیوپاتیک اپیلپسی: به دنبال تب دچار تشدید دفعات تشنج.

تب بالا در بارداری می‌تواند تراتوژن باشد.

تب کمتر از ۳۹ درجه در کودک سالم معمولاً نیاز به درمان ندارد.
\begin{itemize}
\item اگر کودک بی‌قرار شد، درمان می‌کنیم.
\end{itemize}

هیدریشن خوب باید انجام شود.

هایپرپایرکسی (\FeverMath{41<}): حتماً باید تب درمان شود.
\textbf{یادداشت نماد:} در پرانتز، ۴۱ سمت چپ «<» نوشته شده است.

\begin{latin}
Acetaminophen 10--15 mg/kg/dose q4h\\
Ibuprofen in children \(>6\) months old at 5--10 mg/kg/dose q8h
\end{latin}

پاشویه یا پتوهای خنک‌کننده در کنترل تب مؤثر نیستند.

داروهای مکمل و گیاهی شواهدی ندارند.

\clearpage
\section*{پیوست: شش تصویر کامل منبع}
% Preserve the requested "ped-infectioud diseases" image-directory spelling.
\IfFileExists{../../ped-infectioud diseases/fever/images/pic01.jpg}{%
  \def\FeverImageBase{../../ped-infectioud diseases/fever/images/}}{%
  \def\FeverImageBase{pediatrics/ped-infectioud diseases/fever/images/}}
\providecommand{\FeverOriginal}[2]{%
  \clearpage\subsection*{اصل تصویر #1: \lr{\texttt{pic#2.jpg}}}%
  \begin{center}
  \includegraphics[width=\linewidth,height=0.87\textheight,keepaspectratio]{\FeverImageBase pic#2.jpg}
  \end{center}}
\FeverOriginal{۱}{01}
\FeverOriginal{۲}{02}
\FeverOriginal{۳}{03}
\FeverOriginal{۴}{04}
\FeverOriginal{۵}{05}
\FeverOriginal{۶}{06}
\FeverFinish
